\documentclass[11pt]{article}
\usepackage[margin=1in]{geometry}
\usepackage{amsmath,amssymb,amsthm,mathtools,booktabs}
\usepackage[hidelinks]{hyperref}
\usepackage{microtype}
\setlength{\emergencystretch}{2em}
\newtheorem{theorem}{Theorem}[section]
\newtheorem{lemma}[theorem]{Lemma}
\newtheorem{proposition}[theorem]{Proposition}
\newtheorem{corollary}[theorem]{Corollary}
\theoremstyle{remark}\newtheorem{remark}[theorem]{Remark}
\newcommand{\R}{\mathbb R}\newcommand{\C}{\mathbb C}
\newcommand{\dd}{\,\mathrm d}
\newcommand{\abs}[1]{\lvert#1\rvert}
\newcommand{\norm}[1]{\lVert#1\rVert}
\DeclareMathOperator{\Var}{Var}
\DeclareMathOperator{\Res}{Res}
\DeclareMathOperator{\diam}{diam}
\DeclareMathOperator{\supp}{supp}
\title{Quadratic rigidity of shortest positive exponential filters}
\author{}
\date{27 November 2026}
\begin{document}\maketitle
\begin{abstract}
We determine the full width Hessian at every reflection-paired spectrum whose optimal positive filter has strictly positive weights. The targets may have different real parts and arbitrary prescribed paired multiplicities. The quadratic form is an explicit integral of variances, and the atom velocities are obtained by a finite polynomial formula. Its kernel consists exactly of deformations preserving reflection symmetry and the weighted mean imaginary part. Consequently the excess over the mean-frequency lower bound is locally comparable to the squared distance from the reflection locus. This resolves the local Hessian and quadratic-rigidity question posed in the positive-filter chapter of ProveIt, beyond its previously treated coalesced configuration. It does not assert a sign rule for every nonprincipal matrix minor, does not address collisions of distinct labelled targets, and uses neither RH nor GRH. The residue certificate, strip identity, and local continuation method used in the proof are inherited from that chapter; the general second variation, explicit velocity formula, and exact normal kernel are the new deductions.
\end{abstract}

\section{Problem and the theorem}

Let $s_1,\ldots,s_r$ be distinct complex numbers with positive imaginary parts, and let $m_j\ge1$ be integers. Put
\[
 N=\sum_{j=1}^r m_j,\qquad
 \kappa=\frac1N\sum_{j=1}^r m_j\Im s_j.
\]
For a compactly supported probability measure $\mu$ on $\R$, write
\[
 F_\mu(s)=\int_\R e^{st}\dd\mu(t).
\]
The width $D(\boldsymbol s,\boldsymbol m)$ is the infimum of $\diam\supp\mu$ over measures for which $F_\mu$ has a zero of order at least $m_j$ at every $s_j$. All derivatives here are complex derivatives of the transform. Translation preserves the zero conditions. A finite feasible measure always exists: convolve two-point measures having a simple zero at each target, with the requisite repetition. Weak compactness after translation shows that the minimum is attained.

A base configuration is \emph{reflection-paired} if there is an involution $j\mapsto j^*$ such that
\begin{equation}\label{eq:mirror}
 m_{j^*}=m_j,\qquad s_{j^*}^0=\overline{s_j^0}+2i\kappa_0,
 \qquad 0<\Im s_j^0<2\kappa_0.
\end{equation}
Set
\begin{equation}\label{eq:base}
 T=2\kappa_0,\qquad h=\frac{\pi}{\kappa_0},\qquad
 z_j=e^{hs_j^0},\qquad
 P(z)=\prod_j(z-z_j)^{m_j}=\sum_{k=0}^N p_kz^k.
\end{equation}
Throughout the local theorem we assume $p_k>0$ for every $k$. The polynomial is real and $p_0>0$. Define $w_k=p_k/P(1)$.

Consider a real parameter $\varepsilon$ and a velocity $\boldsymbol v\in\C^r$ satisfying the single weighted-mean constraint
\begin{equation}\label{eq:fixedmean}
 \sum_jm_j\Im v_j=0.
\end{equation}
This condition does not require any individual imaginary part to remain fixed. Write
\begin{equation}\label{eq:defectvelocity}
 d_j=\frac{v_j-\overline{v_{j^*}}}{i}.
\end{equation}
Then $d_{j^*}=\overline{d_j}$ and $\sum_jm_jd_j=0$.

\begin{theorem}[Full mirror Hessian]\label{thm:hessian}
Under \eqref{eq:mirror}--\eqref{eq:fixedmean}, the following conclusions hold.
\begin{enumerate}
\item In a neighborhood of the base, the unique minimizing probability measure with left endpoint zero has $N+1$ positive atoms. Its weights and nodes are real analytic in the real and imaginary parts of the targets, with fixed multiplicities. At the base it is
\[
 \mu_0=\sum_{k=0}^Nw_k\delta_{kh},\qquad D_0=Nh.
\]
\item Along the linear path $s_j(\varepsilon)=s_j^0+\varepsilon v_j$, let $u_k$ be the derivative at zero of the $k$th atom position. Then $u_0=u_N=0$ and
\begin{equation}\label{eq:U}
 U(z):=\sum_{k=0}^Nw_ku_kz^k
 =-\frac{h}{T P(1)}\sum_{j=1}^r m_jd_jz_j\frac{P(z)}{z-z_j}.
\end{equation}
The right side is a real polynomial of degree at most $N-1$ with constant coefficient zero.
\item The optimal width has expansion
\begin{equation}\label{eq:expansion}
 D(\boldsymbol s^0+\varepsilon\boldsymbol v,\boldsymbol m)
 =Nh+\varepsilon^2\mathcal Q(\boldsymbol v)+O(\varepsilon^3),
\end{equation}
where
\begin{align}
 \mathcal Q(\boldsymbol v)
 &=\int_{\R}\Var_{\pi_x}(u_k)\dd x,
 &\pi_x(k)&=\frac{p_ke^{khx}}{P(e^{hx})},\label{eq:variance}\\
 &=\frac1h\sum_{0\le k<\ell\le N}p_kp_\ell(u_k-u_\ell)^2
 \int_0^\infty\frac{y^{k+\ell-1}}{P(y)^2}\dd y.\label{eq:gram}
\end{align}
Thus the Hessian restricted to the weighted-mean hyperplane is $2\mathcal Q$.
\item This Hessian is positive semidefinite, and
\begin{equation}\label{eq:kernel}
 \mathcal Q(\boldsymbol v)=0
 \quad\Longleftrightarrow\quad
 v_j=\overline{v_{j^*}}\quad\hbox{for every }j.
\end{equation}
Its kernel is exactly the tangent space to the reflection locus with fixed $\kappa_0$. No extra flat direction occurs anywhere in the strictly positive coefficient locus.
\end{enumerate}
\end{theorem}

All integrals in \eqref{eq:gram} converge: $1\le k+\ell\le2N-1$ and $P(y)$ has positive coefficients. The formula involves only polynomial coefficients, finite arithmetic, and positive one-dimensional integrals. It does not require solving the perturbed optimization problem.

The claim about positive semidefiniteness is stronger when accompanied by \eqref{eq:kernel} and \eqref{eq:U}: nonnegativity alone also follows from the global minimum bound proved below. Signs of nonprincipal minors depend on the coordinates used, and are not claimed here. The strict total positivity proved in the source for special coefficient coordinates at a coalesced frequency is a separate assertion.

\section{The global lower bound for arbitrary real parts}

The following observation allows different real parts in the source's mean-frequency theorem.

\begin{lemma}[Strip identity]\label{lem:strip}
Let $G(s)=\sum_{k=0}^M a_ke^{t_ks}$, where $0=t_0<\cdots<t_M=L$, $a_k>0$, and $\sum a_k=1$. For every $T>0$,
\begin{equation}\label{eq:strip}
 \int_\R\log\frac{G(x)}{\abs{G(x+iT)}}\dd x
 =\frac{LT^2}{2}-2\pi\sum_{0<\Im\rho<T}(T-\Im\rho),
\end{equation}
where the sum is over zeros of $G$, with multiplicity. The integral is nonnegative.
\end{lemma}
\begin{proof}
Endpoint dominance confines the zeros in any horizontal strip to a bounded vertical rectangle. The logarithmic singularities at zeros are integrable. In distributions, $\Delta\log|G|=2\pi\sum_\rho\delta_\rho$. Integrating its $xx$ derivative over the real line gives the difference $L-0$ of the two endpoint slopes. Thus the left side $J(T)$ satisfies $J''(T)=L-2\pi\sum_\rho\delta_{\Im\rho}(T)$ for $T>0$. Also $J(0)=J'(0)=0$, the latter by conjugation symmetry. Integrate twice. Finally $|G(x+iT)|\le G(x)$ by the triangle inequality. Boundary zeros are handled by limits; their weight in the sum is zero.
\end{proof}

\begin{theorem}[Complex-target mean bound]\label{thm:mean}
For arbitrary distinct targets with positive imaginary parts,
\[
 D(\boldsymbol s,\boldsymbol m)\ge\frac{N\pi}{\kappa}.
\]
Equality holds precisely when all target imaginary parts are below $2\kappa$, the multiset is invariant under $s\mapsto\bar s+2i\kappa$, and
\[
 P(z)=\prod_j\bigl(z-e^{\pi s_j/\kappa}\bigr)^{m_j}
\]
has nonnegative coefficients. In that case the unique optimum, after translating its left endpoint to zero, has masses $[z^k]P(z)/P(1)$ at $k\pi/\kappa$.
\end{theorem}
\begin{proof}
Carath\'eodory's theorem applied to the mass and real and imaginary jet constraints replaces an optimal measure by a finite one of no greater width. Translate it to $[0,L]$ and remove zero weights. Apply Lemma~\ref{lem:strip} at $T=2\kappa$. The zero sum is at least
\[
 \sum_jm_j(T-\Im s_j)_+\ge NT-\sum_jm_j\Im s_j=N\kappa.
\]
Since $J(T)\ge0$, this gives $L\ge N\pi/\kappa$.

At equality, $J(T)=0$ and no target lies above height $T$. The integrand vanishes almost everywhere and hence everywhere away from its possible singularities; a singularity would itself force a positive integral. Equality in the triangle inequality at $x=0$ forces every atom of the finite replacement into $h\mathbb Z$, $h=\pi/\kappa$. Then $G(x+iT)=G(x)>0$ for real $x$, so no target has imaginary part $T$. The transform is $A(e^{hs})$ for a polynomial $A$ of degree $N$ with nonnegative coefficients. The exponential map is injective on the target strip and locally invertible, so $A=P/P(1)$. Real coefficients give reflection pairing. Conversely, the stated coefficients produce an admissible measure of width $Nh$. To conclude uniqueness among all measures, not just finite replacements, use the certificate in Lemma~\ref{lem:certificate}, which does not require strict coefficient positivity. Any optimal measure integrates it to zero and is therefore supported on the same grid; polynomial divisibility then fixes all its masses.
\end{proof}

\section{Why the minimizing measure is analytic}

We give the local argument, including its global optimality certificate, so that Theorem~\ref{thm:hessian} does not rest on numerical stationarity.

\begin{lemma}[Strict certificate at the base]\label{lem:certificate}
At a reflection-paired base as in \eqref{eq:mirror}, there is a real function in the span of the real and imaginary parts of $t^ke^{s_j^0t}$, $0\le k<m_j$, that is positive on $(0,Nh)$ except at $h,2h,\ldots,(N-1)h$. Its interior zeros are double and its endpoint zeros are simple, with inward positive slopes.
\end{lemma}
\begin{proof}
Put $q_j=-z_j$, $Q(z)=\prod_j(z-q_j)^{m_j}$ and, using the principal logarithm,
\[
 E(x)=\sum_j\Res_{z=q_j}\frac{z^{x-1}}{Q(z)},\qquad
 g(hx)=(-1)^{N+1}\sin(\pi x)E(x).
\]
The $q_j$ avoid the negative real axis and occur in conjugate pairs. For $u\ge0$ we have $Q(-u)=(-1)^N|Q(-u)|$. A keyhole contour gives, for $0<x<N$,
\begin{equation}\label{eq:certificate}
 g(hx)=\frac{\sin^2(\pi x)}{\pi}
 \int_0^\infty\frac{u^{x-1}}{|Q(-u)|}\dd u.
\end{equation}
The inner and outer circular contour contributions tend to zero since $0<x<N$. The integral is positive, behaves as $1/(|Q(0)|x)+O(1)$ at zero and as $1/(N-x)+O(1)$ at $N$. These give all strict signs and orders claimed. Residue expansion expresses $E$ as a sum of polynomials times $e^{x\log q_j}$ of the prescribed degrees. Multiplication by $\sin\pi x$ produces exponents $s_j^0$ and $s_j^0-iT=\overline{s_{j^*}^0}$, proving membership in the required span.
\end{proof}

\begin{proposition}[Local uniqueness and continuation]\label{prop:analytic}
Part (1) of Theorem~\ref{thm:hessian} holds.
\end{proposition}
\begin{proof}
Use the $2N+1$ unknowns consisting of $N+1$ weights and $N$ nodes, with the first node anchored at zero. The equations are mass one and the $2N$ real jet constraints. At the base, a variation of the transform has the form
\[
 \delta F(s)=A(e^{hs})+sB(e^{hs}),\qquad A(1)=0,\quad B(0)=0,
\]
where $A,B$ are real polynomials of degree at most $N$. If this variation is in the equation Jacobian's kernel, all target jets vanish. Conjugating the reflected partner equations shows that the same jets of $\delta F(s-iT)$ vanish at $s_j^0$. Since $e^{-ihT}=1$, subtraction gives $iTB(e^{hs})$; hence $P$ divides $B$. Its degree and $B(0)=0\ne P(0)$ imply $B=0$. Now $P$ divides $A$ and $A(1)=0\ne P(1)$, so $A=0$. Positivity of all base weights implies that every node variation vanishes. The real Jacobian is invertible, and the analytic implicit function theorem gives an analytic solution with positive weights and ordered nodes.

Let $Y(t)$ be the vector of real moment functions. The moment columns $Y(t_k)$ and the interior derivative columns $Y'(t_k)$ span a hyperplane in $\R^{2N}$: the weight relation supplies one dependence, and adjoining the last derivative column gives full rank by the preceding Jacobian calculation. Choose its analytic normal vector so that its scalar product with the last derivative column is $-1$. Its scalar product with $Y(t)$ is an analytic family of certificates vanishing at all nodes, with derivative zero at every interior node. At the base it is a positive multiple of Lemma~\ref{lem:certificate}'s certificate. The strict endpoint slopes, positive second derivatives at interior nodes, and positive minimum away from their neighborhoods persist under small perturbations. Therefore each certificate is nonnegative on the whole proposed support, with precisely its nodes as zeros.

Any feasible measure of no greater width, translated to start at zero, must integrate this certificate to zero and so is supported on these nodes. The mass and moment columns have full column rank by the same Jacobian calculation, so its weights are the proposed weights. The positive endpoint weights force the full width. Thus the analytic solution is the unique global optimum up to translation.
\end{proof}

\section{Exact defect identity and second variation}

Write the analytic optimum as $G_\varepsilon(s)=\sum_{k=0}^Nw_k(\varepsilon)e^{t_k(\varepsilon)s}$, with $t_0=0$. Along a fixed-mean path, the strip height $T=2\kappa_0$ is fixed.

\begin{lemma}[Exact local defect]\label{lem:defect}
For all sufficiently small real $\varepsilon$,
\begin{equation}\label{eq:exactdefect}
 D(\varepsilon)-Nh
 =\frac{2}{T^2}\int_\R
 -\log\left|\sum_{k=0}^N\pi_x(k,\varepsilon)e^{iTt_k(\varepsilon)}\right|\dd x,
 \quad
 \pi_x(k,\varepsilon)=\frac{w_k(\varepsilon)e^{xt_k(\varepsilon)}}{G_\varepsilon(x)}.
\end{equation}
\end{lemma}
\begin{proof}
Uniform lower bounds on endpoint weights and adjacent node gaps confine every zero of $G_\varepsilon$ in $0\le\Im s\le T$ to a common vertical rectangle: on either real tail the endpoint term dominates the sum of all other terms. The base transform $P(e^{hs})/P(1)$ has exactly $N$ zeros in $0<\Im s<T$, namely the prescribed targets with their multiplicities, and no zeros on its horizontal boundary. Rouch\'e's theorem on a large fixed rectangle therefore gives exactly $N$ zeros of $G_\varepsilon$ in the strip. The imposed target jets already account for all of them. Their weighted sum of imaginary parts remains $N\kappa_0$. Lemma~\ref{lem:strip} gives $J(T)=T^2(D-Nh)/2$. Dividing by $G_\varepsilon(x)$ gives \eqref{eq:exactdefect}.
\end{proof}

\begin{proof}[Proof of Theorem~\ref{thm:hessian}, parts (2)--(4)]
The mean bound and the fixed-mean condition make $D(\varepsilon)\ge Nh$ with equality at zero. Analyticity gives $D'(0)=0$, hence $u_N=0$; anchoring gives $u_0=0$.

For a probability vector $a$ and a real vector $u$,
\[
 -\log\left|\sum_ka_ke^{iT\varepsilon u_k}\right|
 =\frac{T^2\varepsilon^2}{2}\Var_a(u)+O(\varepsilon^3).
\]
At the base all phases $e^{iTkh}$ equal one. Derivatives of the weights and of the tilt probabilities do not enter the quadratic coefficient, because the expression is identically zero when all phases agree. The same is true of second node derivatives. Differentiation under the integral in \eqref{eq:exactdefect} is valid: on each tail one endpoint carries all but $O(e^{-c|x|})$ of the probability mass, uniformly in small $\varepsilon$, and factoring out its phase bounds the first three parameter derivatives by $C(1+|x|)^3e^{-c|x|}$. On the central compact interval the expression inside the modulus is nonzero near the base. This proves \eqref{eq:variance} and \eqref{eq:expansion}. Expand variance as $\sum_{k<\ell}\pi_k\pi_\ell(u_k-u_\ell)^2$ and substitute $y=e^{hx}$ to obtain \eqref{eq:gram}.

It remains to derive the promised atom velocities. Put $U(z)=\sum_kw_ku_kz^k$ and $A(z)=\sum_kw'_k(0)z^k$. The first transform variation is $\delta F(s)=A(e^{hs})+sU(e^{hs})$. Differentiating the zero of order $m_j$ at the moving target shows that
\[
 \delta F(s)+v_jF_0'(s)
\]
vanishes to order at least $m_j$ at $s_j^0$. Conjugate the reflected partner equation and shift by $iT$. Subtraction shows that
\[
 T U(z)+\frac{h d_j}{P(1)}zP'(z)
 \quad\hbox{is divisible by }(z-z_j)^{m_j}.
\]
The polynomial on the right of \eqref{eq:U} satisfies these congruences: modulo $(z-z_j)^{m_j}$, the only nonzero term in the sum is $m_jd_jz_jP(z)/(z-z_j)$, while $zP'(z)$ is congruent to $m_jz_jP(z)/(z-z_j)$. It has degree at most $N-1$, real coefficients by conjugation pairing, and constant coefficient proportional to $\sum_jm_jd_j=0$. The actual $U$ also has degree at most $N-1$ since $u_N=0$. Their difference is divisible by $P$ and has smaller degree, proving \eqref{eq:U}.

Every $\pi_x(k)$ is positive. The integral of variance vanishes precisely when all $u_k$ agree; anchoring then makes all of them zero. Formula \eqref{eq:U}, divided by $P$, gives the rational identity
\[
 U=0\quad\Longleftrightarrow\quad
 \sum_j\frac{m_jd_jz_j}{z-z_j}=0.
\]
Distinct poles and nonzero $z_j$ force each $d_j=0$. This is exactly \eqref{eq:kernel}.
\end{proof}

\section{Quantitative rigidity and changing the mean}

The normal quadratic form comes with elementary explicit bounds. Write $w_{\min}=\min_kw_k$. For $N\ge2$,
\begin{equation}\label{eq:coercivity}
 \frac{2e^{-2N}w_0w_{\min}}{h}\sum_{k=1}^{N-1}u_k^2
 \le\mathcal Q(\boldsymbol v)
 \le\frac1h\sum_{k=1}^{N-1}w_k
 \left(\frac1{w_0k}+\frac1{w_N(N-k)}\right)u_k^2.
\end{equation}
For the lower bound use $\pi_x(j)\ge e^{-N}w_j$ on $|hx|\le1$ and keep pairs $(0,k)$ in the variance. For the upper bound use $\Var(u)\le\sum_k\pi_x(k)u_k^2$, and dominate $\pi_x(k)$ by $(w_k/w_N)e^{-(N-k)hx}$ on $x\ge0$ and by $(w_k/w_0)e^{khx}$ on $x\le0$. Formula \eqref{eq:U} gives the remaining finite linear map from spectral defects to node velocities, so singular values of that explicit map give computable spectral coercivity constants.

\begin{corollary}[Local squared-distance law]\label{cor:rigidity}
Fix a strictly positive base and fixed target labels and multiplicities. For a nearby spectrum set $\kappa=N^{-1}\sum m_j\Im s_j$ and
\[
 \mathcal E(\boldsymbol s)^2
 =\sum_jm_j|s_j-\overline{s_{j^*}}-2i\kappa|^2.
\]
There are constants $0<c<C<\infty$ and a neighborhood of the base such that
\[
 c\mathcal E(\boldsymbol s)^2
 \le D(\boldsymbol s,\boldsymbol m)-\frac{N\pi}{\kappa}
 \le C\mathcal E(\boldsymbol s)^2.
\]
\end{corollary}
\begin{proof}
The linear projection
\[
 s_j^{\rm mir}=\frac{s_j+\overline{s_{j^*}}+2i\kappa}{2}
\]
preserves the weighted mean and takes spectra to the reflection locus. It remains within the strictly positive coefficient locus near the base. The difference $\boldsymbol s-\boldsymbol s^{\rm mir}$ is a normal displacement whose Euclidean norm, with multiplicity weights, is $\mathcal E/2$. The width excess vanishes identically on this reflection locus. Theorem~\ref{thm:hessian} gives a positive definite Hessian on its normal space, uniformly on a sufficiently small closed neighborhood. Taylor expansion about $\boldsymbol s^{\rm mir}$, with its uniform cubic remainder, proves both inequalities. If the normal space is zero-dimensional, the excess and defect both vanish and the assertion is vacuous.
\end{proof}

For a general linear velocity, put $a=N^{-1}\sum m_j\Im v_j$ and $\widehat v_j=v_j-(a/\kappa_0)s_j^0$. This adjusted velocity has zero weighted mean. Positive real scaling of every target divides the optimal width by the same scale, so Theorem~\ref{thm:hessian} yields
\begin{equation}\label{eq:allhessian}
 D(\boldsymbol s^0+\varepsilon\boldsymbol v)
 =Nh-\frac{Nh\,a}{\kappa_0}\varepsilon
 +\left(\frac{Nh\,a^2}{\kappa_0^2}+\mathcal Q(\widehat{\boldsymbol v})\right)\varepsilon^2
 +O(\varepsilon^3).
\end{equation}
The possibly quadratic term in the rescaled target path does not change the excess coefficient, because the gradient of the excess vanishes on the reflection locus. Thus the full Hessian is also positive semidefinite. Its nullspace consists exactly of fixed-mean reflection-preserving velocities. This is a pointwise Hessian assertion on the reflection locus, not a claim that the width is convex throughout a neighborhood.

\section{An explicit three-mode formula}

Take $\boldsymbol s^0=(i,2i,3i)$, all multiplicities one. Then $\kappa_0=2$, $h=\pi/2$, $P(z)=1+z+z^2+z^3$, and $w_k=1/4$. Write $v_j=a_j+ib_j$ with $b_1+b_2+b_3=0$, and put
\[
 X=b_1+b_3=-b_2,\qquad Y=a_1-a_3.
\]
Direct substitution in \eqref{eq:U} gives
\[
 u_1=\frac\pi4(X-Y),\qquad
 u_2=-\frac\pi4(X+Y),\qquad u_0=u_3=0.
\]
The integrals needed for \eqref{eq:gram} are
\[
 \int_0^\infty\frac{\dd y}{P(y)^2}=\frac\pi8,
 \quad\int_0^\infty\frac{y\dd y}{P(y)^2}=\frac\pi8-\frac14,
 \quad\int_0^\infty\frac{y^2\dd y}{P(y)^2}=\frac12-\frac\pi8;
\]
the remaining integrals follow by $y\mapsto1/y$. Therefore
\begin{equation}\label{eq:triple}
 \boxed{\quad D(\boldsymbol s^0+\varepsilon\boldsymbol v)
 =\frac{3\pi}{2}+\frac\pi{16}
 \bigl(3X^2+(\pi-1)Y^2\bigr)\varepsilon^2+O(\varepsilon^3).\quad}
\end{equation}
Frequency asymmetry and unequal real parts give separate strictly positive costs. Common changes of all real parts, changes of the middle real part, and opposite shifts of the outer frequencies preserve the mirror condition and have no quadratic cost; in fact they preserve the exact mean-bound width while coefficient positivity persists.

\section{Source comparison, scope, and verification}

The source is the positive-filter chapter at ProveIt commit \texttt{ce35f320bc88} \cite{ProveIt}. The source labels \texttt{fil:lem:strip}, \texttt{fil:thm:reflect}, and \texttt{fil:thm:perturb} supply the strip identity, residue certificate, and local analytic optimum. The result labelled \texttt{fil:thm:differential} gives the first derivative. The source's \texttt{fil:thm:hessian} treats the coalesced configuration in symmetric-polynomial coordinates. Open item 2 in \texttt{fil:sec:open} asks for the Hessian at general positive mirror configurations. The new part here is the closed formula \eqref{eq:U} at an arbitrary such configuration, the resulting full complex-target and multiplicity-aware Hessian, the exact kernel, and the local squared-distance law. We have retained proofs of the inherited mechanisms and identified them explicitly. The complex-target extension of the mean bound is a short corollary of the same inherited strip identity.

This is an unconditional finite-dimensional extremal theorem. In the repository's arithmetic applications, a positive logarithmic average multiplies a zero contribution by $F_\mu(s)$; the theorem gives the exact quadratic increase in the smallest averaging range when prescribed exponents depart from a mirror configuration. It does not prove a new statement about locations or spacings of zeros of $L$-functions. It does not settle the global three-frequency phase diagram, the optimal logarithmic-support second term, or a prominent named number-theoretic conjecture.

The accompanying script checks the exact three-mode polynomial and integral identities and compares the predicted coefficient with high-precision solutions of the full mass and jet equations, including a repeated-target example. These numerical checks illustrate the theorem; the proof does not depend on them. No Lean or Rocq formalization is claimed. The comparison establishes an extension of the inspected source, not publication priority over all mathematical literature. The strip inequality has classical antecedents in inequalities for zeros of polynomials with nonnegative coefficients \cite{EF}.

\begin{thebibliography}{9}
\bibitem{ProveIt}
V.~Reshetnikov and contributors, \emph{ProveIt: Generalized Riemann Hypothesis, Article, Positive filters that cancel zeros}, inspected commit \texttt{ce35f320bc88} (full identifier in the URL).
\url{https://github.com/VladimirReshetnikov/ProveIt/blob/ce35f320bc88c02093e479f5fe5003ff86b2795f/Analysis/GeneralizedRiemannHypothesis/Article/sections/fil.tex}.
\bibitem{EF}
A.~Eremenko and A.~Fryntov, \emph{Remarks on the Obrechkoff inequality}, arXiv:1409.4849.
\url{https://arxiv.org/abs/1409.4849}.
\end{thebibliography}
\end{document}
